\originalchapter{Jeans 引力不稳定与旋转盘的稳定窗口}

本章对应 8.902 第7讲. 定常分布函数只说明一组轨道可以不变, 并未保证小扰动不会增长. 核心问题是随机运动或压力来得及抹平密度起伏吗, 旋转的回复作用又保护哪些波长? 先完整推导流体近似的色散关系, 再区分恒星盘的无碰撞结果. 原课以尺度估计引入 Toomre 条件, 其中的近似系数与严格动力学系数不能混为一谈.

\originalsection{均匀背景的近似及闭合假设}
对理想可压缩流体, 取 $\rho=\rho_0+\delta\rho$, 平均速度 $\boldsymbol u=\delta\boldsymbol u$, 线性压力扰动 $\delta P=c_s^2\delta\rho$. $c_s^2$ 是沿所采用热力学过程的导数 $\partial P/\partial\rho$, 绝热与等温值未必相同. 用流体的连续性与 Euler 方程,
\[
\partial_t\delta\rho+\rho_0\nabla\cdot\delta\boldsymbol u=0,
\qquad
\partial_t\delta\boldsymbol u=-\nabla\delta\Phi-
\frac{c_s^2}{\rho_0}\nabla\delta\rho,
\quad \nabla^2\delta\Phi=4\pi G\delta\rho.
\]
严格静态、无限均匀且非零密度的背景不能同时有 $\nabla\Phi_0=0$ 与完整 Poisson 方程. 上式只给扰动密度求引力, 将均匀背景力另行抵消, 是所谓 Jeans swindle. 它可以作为局部近似或适当周期模型, 不是已建立的无限孤立平衡宇宙. 后面在膨胀背景中写扰动方程时, 背景引力将通过 Friedmann 方程处理.

对无碰撞恒星, Jeans 矩方程中的应力尚无普适状态方程. 将其简化为同一等向 $\delta P=c_s^2\delta\rho$ 是额外流体闭合, 只能给速度弥散支持的尺度直觉; 完整动力学需解分布函数的线性响应, 还可能有相混合效应.

\originalsection{波长与增长率从色散关系读出}
对连续性再微分并用另外两式消去速度和势,
\[
\partial_t^2\delta\rho-c_s^2\nabla^2\delta\rho
-4\pi G\rho_0\delta\rho=0.
\]
代入 $\delta\rho\propto\exp[\ii(\boldsymbol k\cdot\boldsymbol x-\omega t)]$,
\[
\omega^2=c_s^2k^2-4\pi G\rho_0.
\]
高 $k$ 短波由压力支持, 低 $k$ 长波有增长模与衰减模. 令
\[
k_J=\frac{\sqrt{4\pi G\rho_0}}{c_s},\qquad
\lambda_J=\frac{2\pi}{k_J}=c_s\sqrt{\frac{\pi}{G\rho_0}},
\]
则 $\lambda>\lambda_J$ 时增长率 $s=\sqrt{4\pi G\rho_0-c_s^2k^2}$. 正好临界时 $\omega=0$, 通常解有常数项及线性时间项, 不能把它当作已有正频率稳定振动.

Jeans 质量的几何约定需说明. 本卷按直径 $\lambda_J$ 的球定义
\[
M_J=\frac{4\pi}{3}\rho_0\left(\frac{\lambda_J}{2}\right)^3
=\frac{\pi^{5/2}}6\frac{c_s^3}{G^{3/2}\rho_0^{1/2}}.
\]
原课把 $\lambda_J$ 当球半径写 $4\pi\rho_0\lambda_J^3/3$, 比本卷大8倍; 波长阈值相同, 质量只因定义不同. 比较两个文献的 $M_J$ 前须核对这个选择. 尺度上 $t_{\rm cross}\sim\lambda/c_s$, $t_{\rm ff}\sim1/\sqrt{G\rho_0}$, 长尺度在压力通信前发生塌缩, 这解释了上述符号, 不用时间尺度估计冒充所有 $\pi$ 因子的精确推导.

\originalsection{薄盘自引力的二维 Fourier 核}
在局部盘面取径向坐标 $x$, 对薄盘扰动 $\delta\Sigma\exp(\ii kx)$, 三维密度为 $\delta\Sigma\delta(z)\exp(\ii kx)$. 离盘面 $\delta\Phi\propto\exp(-|k||z|)$, Poisson 方程积分跨盘面给导数跳跃, 因而
\[
\delta\Phi_k(z)=-\frac{2\pi G}{|k|}\delta\Sigma_k\ee^{-|k||z|}.
\]
负号与引力吸引一致. 由此盘的引力项正比于 $|k|$, 与均匀三维介质的常数 $4\pi G\rho_0$ 不同. 这就是不能直接把三维 Jeans 长度换一个密度就当作严格盘判据的原因.

\originalsection{旋转气体盘的完整局部色散关系}
在导引中心旋转的局部参考系中, 设径向扰动速度为 $u$, 切向为 $v$, 只取轴对称、短波和线性扰动, 背景面密度 $\Sigma_0$ 近似常数. 含 Coriolis 与背景剪切的方程为
\begin{align*}
-\ii\omega\delta\Sigma+\ii k\Sigma_0u&=0,\\
-\ii\omega u-2\Omega v&=-\ii k\left(c_s^2\frac{\delta\Sigma}{\Sigma_0}+\delta\Phi\right),\\
-\ii\omega v+\frac{\kappa^2}{2\Omega}u&=0.
\end{align*}
先由第一式 $\delta\Sigma=k\Sigma_0u/\omega$, 第三式 $v=\kappa^2u/(2\ii\omega\Omega)$, 再代入径向式及薄盘势核, 得
\[
\omega^2=\kappa^2+c_s^2k^2-2\pi G\Sigma_0|k|.
\]
转动回复项 $\kappa^2$ 保护足够长的波, 压力项 $c_s^2k^2$ 保护短波, 中间尺度可能不稳定. 将 $k\ge0$ 时右端配方,
\[
\omega^2=c_s^2\left(k-\frac{\pi G\Sigma_0}{c_s^2}\right)^2
+\kappa^2-\frac{\pi^2G^2\Sigma_0^2}{c_s^2}.
\]
故全部局部轴对称波无指数增长的条件为
\[
Q_g\equiv\frac{c_s\kappa}{\pi G\Sigma_0}\ge1.
\]
若 $Q_g<1$, 不稳定波数区间由二次方程两根 $k_-<k<k_+$ 给出. 冷盘 $c_s=0$ 时只有低 $k$ 转动保护, 更短波仍不稳定. 随机运动与转动必须共同作用才可能封闭不稳定窗口.

\begin{center}
\begin{tikzpicture}
\begin{axis}[width=8.2cm,height=5cm,xlabel={$q=kc_s/\kappa$},ylabel={$\omega^2/\kappa^2$},xmin=0,xmax=3,ymin=-1,ymax=5,grid=major,legend style={font=\small}]
\addplot[structurecolor,thick,domain=0:3]{1+x^2-2*x/0.8};\addlegendentry{$Q_g=0.8$}
\addplot[black,dashed,domain=0:3]{1+x^2-2*x/1.2};\addlegendentry{$Q_g=1.2$}
\addplot[gray,domain=0:3]{0};
\end{axis}
\end{tikzpicture}
\end{center}
图按色散式直接绘制, 负的频率平方才是增长区, 不是频率为负. 有限厚度减弱小尺度引力响应; 多组分气体与恒星共同耦合, 又会改变单组分判据. $Q_g>1$ 只控制这里的局部轴对称线性模, 不保证没有非轴对称棒、螺旋或非线性演化.

\originalsection{恒星盘的系数来自分布函数响应}
无碰撞 Gaussian 径向速度分布的薄盘结果通常写为
\[
Q_*\equiv\frac{\sigma_R\kappa}{3.36G\Sigma_0}\ge1.
\]
这里以 Toomre 1964 原论文的临界结果作引用输入, 完整速度空间线性响应不由前面流体闭合证明. 在零频率边界, 其径向随机运动减弱因子含 $[1-\ee^{-\chi}I_0(\chi)]/\chi$, $\chi=(k\sigma_R/\kappa)^2$, $I_0$ 为修正 Bessel 函数. 临界常数可独立重算为
\[
C=2\pi\max_{q>0}\frac{1-\ee^{-q^2}I_0(q^2)}q
\simeq3.35828,
\quad q_{\max}\simeq0.97370.
\]
课程尺度估计 $\sqrt{32/3}\simeq3.266$ 接近这个数, 但其后将3.26写成 Toomre 精确计算的系数并不准确. 本卷区分近似估算、流体的 $\pi$ 与恒星 Gaussian 响应的3.36. 引用原论文只用于此事实核对及常数复算, 不收录论文原件或受限图; 详见来源补充清单.

从这一章回看前面两章, $\kappa$ 来自单星近圆轨道的回复作用, $\sigma_R$ 来自系综的随机运动, $\Sigma_0$ 决定集体自引力. 三者在稳定性中才相遇, 不能只从维里定理知道速度尺度就自动断言稳定. 后面的恒星形成需进一步加入冷却与反馈, 稳定性阈值只是其中的一环.
